\begin{figure}[!t]
  \centering
  \begin{tikzpicture}[
      font=\scriptsize,
      box/.style={draw=figdata!75, fill=white, rounded corners=2pt, align=left, inner sep=3pt,
                  execute at begin node=\hyphenpenalty 10000\relax},
      ex/.style={box, text width=6.85cm, anchor=north west},
      item/.style={align=left, inner sep=0pt, anchor=north west, text width=5.8cm,
                   execute at begin node=\hyphenpenalty 10000\relax},
      lane/.style={rounded corners=3pt, inner sep=4pt},
      dlane/.style={lane, draw=figdata!60, fill=figdata!5},
      plane/.style={lane, draw=figprofile!60, fill=figprofile!7},
      heading/.style={font=\scriptsize\bfseries, inner sep=0pt, anchor=south west},
      meta/.style={align=left, inner sep=0pt, anchor=north west, text width=7.06cm,
                   execute at begin node=\hyphenpenalty 10000\relax},
      btitle/.style={inner sep=0pt, anchor=north west},
      plink/.style={-{Stealth[length=4pt]}, draw=figprofile, line width=0.5pt},
      pnote/.style={font=\scriptsize, inner sep=0pt, text=figprofile!85!black, align=center},
    ]
    \def\xr{8.28}

    \node[meta] (hdr) at (0,0)
      {\textbf{Audiência de 16/08/2023 sobre:} Debate sobre a interpretação do artigo 142 da
       Constituição [\ldots]\enspace\textcolor{black!55}{$\cdot$\enspace turno 30}};

    \node[ex] (l1) at ($(hdr.south west)+(0,-4pt)$)
      {``[\ldots] se está sendo sentida a falta do poder moderador no Brasil hoje, é porque ele pode
       estar sendo exercido de forma errada. Nós precisamos, sim, de um poder moderador, mas não
       podemos nos esquecer nunca das palavras do Prof.~Ives Gandra [\ldots]''};
    \node[ex] (l2) at ($(l1.south west)+(0,-3pt)$)
      {``Constituição não se relativiza. Ela é o que nela está escrito, e o que está escrito tem que
       valer. [\ldots] O Constituinte, quando escreveu o art.~53, foi sábio: colocou a palavra
       `quaisquer' [\ldots] eu invoco minha inviolabilidade [\ldots]''};
    \node[ex] (l3) at ($(l2.south west)+(0,-3pt)$)
      {``Eu sou Deputado Federal pelo Rio Grande do Norte [\ldots]''};
    \node[ex] (l4) at ($(l3.south west)+(0,-3pt)$)
      {``Fui General, Comandante de Brigada, e sempre fui contra a violência. [\ldots] Minha passagem
       para a reserva se deu exatamente porque eu não concordava com uma ordem absurda.''};

    \node[item, text=black!55] (full) at (\xr,0) {perfil completo: 621 palavras em 29 itens};
    \node[btitle] (t1) at (\xr,0 |- l1.north) {{\color{figprofile!85!black}\textbf{Posições}}\enspace\textcolor{black!55}{P1 a P16}};
    \node[item] (r1) at ($(t1.south west)+(0,-1.5pt)$)
      {\textbf{P1.} Defende a existência de um poder moderador no Brasil, afirmando que, se há a sensação de sua
       falta, é porque ele pode estar sendo exercido de forma errada.};
    \path let \p1=(r1.south), \p2=(l2.north) in coordinate (y2) at (\xr,{min(\y1-4pt,\y2)});
    \node[btitle] (t2) at (y2) {{\color{figprofile!85!black}\textbf{Critérios e valores}}\enspace\textcolor{black!55}{C1 a C4}};
    \node[item] (r2) at ($(t2.south west)+(0,-1.5pt)$)
      {\textbf{C1.} Baseia-se na literalidade da Constituição para defender que ela não seja relativizada e para
       invocar a inviolabilidade do art.~53.};
    \path let \p1=(r2.south), \p2=(l3.north) in coordinate (y3) at (\xr,{min(\y1-4pt,\y2)});
    \node[btitle] (t3) at (y3) {{\color{figprofile!85!black}\textbf{Alinhamentos declarados}}\enspace\textcolor{black!55}{A1 a A4}};
    \node[item] (r3) at ($(t3.south west)+(0,-1.5pt)$)
      {\textbf{A1.} Declara-se Deputado Federal pelo Rio Grande do Norte.};
    \path let \p1=(r3.south), \p2=(l4.north) in coordinate (y4) at (\xr,{min(\y1-4pt,\y2)});
    \node[btitle] (t4) at (y4) {{\color{figprofile!85!black}\textbf{Forma de argumentar}}\enspace\textcolor{black!55}{F1 a F5}};
    \node[item] (r4) at ($(t4.south west)+(0,-1.5pt)$)
      {\textbf{F1.} Utiliza a própria trajetória profissional e experiência de vida (carreira militar e segurança
       pública) para embasar suas opiniões.};

    \begin{scope}[on background layer]
      \node[dlane, fit=(hdr)(l1)(l4)(hdr.west |- r4.south)] (bfal) {};
      \node[plane, fit=(full)(t1)(r1)(r4)(r4.west |- l4.south)] (bper) {};
    \end{scope}
    \node[heading, text=figdata!90!black] at ($(bfal.north west)+(0,2pt)$) {Falas de treino};
    \node[heading, text=figprofile!85!black] at ($(bper.north west)+(0,2pt)$)
      {Perfil de GENERAL GIRÃO};

    \foreach \i in {1,...,4}
      \draw[plink] (l\i.east) -- ($(r\i.west)+(-2pt,0)$);
    \node[pnote, anchor=north] at ($(bfal.north east)!0.5!(bper.north west)+(0,-3pt)$)
      {LLM\\gerador};
  \end{tikzpicture}
  \caption{Primeiros itens dos quatro blocos do perfil que o Gemma 4 31B escreveu para GENERAL
  GIRÃO, reproduzidos sem correção e com os identificadores que a simulação atribui aos itens, ao lado dos
  trechos das falas de treino que os sustentam. O ator foi
  escolhido por regra fixa: entre os atores com pergunta de teste e com perfil de quatro blocos,
  excluídos os dois do teste preliminar (Seção~\ref{sec:metodo-simulacao}), o de número mediano de
  audiências de treino e, no empate, aquele cujo total de palavras nos turnos de treino, incluídos os
  de presidência, fica mais próximo da mediana do grupo empatado. Os quatro trechos mostrados vêm do
  mesmo turno, o mais longo das duas audiências de treino do ator, com 1.001 das suas 1.745 palavras.
  Os cortes nos trechos e no assunto estão marcados com {[\ldots]}.}
  \label{fig:perfil}
\end{figure}
